%%%%%%%%%%%%%%%%%%%%%%%%%%%%%%%%%
\section{Kết luận}

\subsection{Đánh giá kết quả nghiên cứu}

Đề tài đặt ra hai câu hỏi nghiên cứu chính và cả hai đã được trả lời thỏa đáng bằng bằng chứng thống kê rõ ràng:

\textbf{Câu hỏi 1 --- Thông số kỹ thuật nào ảnh hưởng đến giá GPU?} Mô hình Log-Linear ($R^2 = 78.54\%$) xác nhận cả 6 thông số đều có ý nghĩa ($p < 0.001$). TDP ảnh hưởng mạnh nhất ($+57.5\%$/SD), tiếp theo là dung lượng RAM, xung nhịp, bus width. Đặc biệt, \texttt{memory\_bus} trở nên có ý nghĩa sau biến đổi log (vốn không có ý nghĩa trong OLS).

\textbf{Câu hỏi 2 --- Chênh lệch giá NVIDIA vs AMD?} Welch's ANOVA ($F = 18.39$, $p = 2.409 \times 10^{-5}$): chênh lệch trung bình \textbf{\$245.43}. Mô hình hồi quy cho thấy cùng thông số, NVIDIA đắt hơn AMD \textbf{45.2\%} --- sự chênh lệch do chiến lược định giá chứ không chỉ do cấu hình.

Kết quả \textbf{vượt kỳ vọng}: $R^2 = 78.54\%$ (mục tiêu $\geq 70\%$). Bước chuyển từ OLS ($R^2 = 29.8\%$) sang Log-Linear là cần thiết do dữ liệu giá GPU có phân phối lệch phải cực đoan (Skewness $= 16.84$).

\subsection{Hạn chế và đề xuất cải tiến}

\begin{enumerate}
  \item \textbf{Phương sai sai số chưa đồng đều} ($BP = 99.04$, $p < 2.2 \times 10^{-16}$). \textit{Đề xuất:} WLS hoặc Robust SE (\texttt{vcovHC()}).
  \item \textbf{Phần dư chưa hoàn toàn chuẩn} ($W = 0.948$, $p = 4.056 \times 10^{-13}$). \textit{Đề xuất:} Biến đổi Box-Cox; CLT với $n = 556$ vẫn đảm bảo suy diễn hợp lệ.
  \item \textbf{Selection bias} do loại 83.7\% quan sát thiếu giá (MNAR). \textit{Đề xuất:} Bổ sung dữ liệu từ Amazon/Newegg hoặc Multiple Imputation.
  \item \textbf{Không xét tương tác biến.} \textit{Đề xuất:} Thêm interaction terms (\texttt{tdp:manufacturer}) và Ramsey RESET test.
\end{enumerate}

\subsection{Ý nghĩa thực tiễn}

\begin{enumerate}
  \item \textbf{Công cụ định giá:} Phương trình Log-Linear cho phép ước lượng giá hợp lý dựa trên thông số kỹ thuật --- nếu giá thực vượt xa giá mô hình, GPU có thể đang bị overpriced.
  \item \textbf{Định lượng ``phí thương hiệu'':} NVIDIA đắt hơn AMD $\approx 45.2\%$ cùng thông số --- cơ sở để người mua đánh giá mức premium.
  \item \textbf{Ngân sách doanh nghiệp:} Mô hình giúp so sánh giá trị thực giữa các GPU cho hệ thống AI/render farm.
  \item \textbf{Tín hiệu thị trường:} GPU ra sau rẻ hơn $\approx 3.8\%$/năm ($\beta_{year} = -0.039$), phản ánh xu hướng Moore's Law.
\end{enumerate}

%%%%%%%%%%%%%%%%%%%%%%%%%%%%%%%%%